% !TEX root = ../main.tex

\chapter{Analysis}
\label{ch:analysis}

\section{Classification of binary32 inputs}
\label{sec:input-classification}
Not every input needs to go through the full pipeline: some have an obvious
result, and some do not represent numbers at all. These inputs can be
identified by examining only the exponent and fraction fields
(\cref{tab:special}), before any computation takes place:

\begin{enumerate}[label=\textbf{\arabic*.}, ref=\arabic*, leftmargin=*]
  \item \textbf{Zero} (exponent and fraction both 0) is passed directly to
    the output, with its sign preserved, since the standard defines both
    $+0$ and $-0$~\cite{ieee754}.
  \item \textbf{Subnormal numbers} (exponent 0, fraction $\neq 0$) are
    smaller than $2^{-12}$, so $\sin x = x$ within the binary32
    precision, and they are passed directly to the output.
  \item \textbf{Tiny normal numbers}, with an exponent field smaller than 115
    ($|x| < 2^{-12}$), are passed directly to the output for the same
    reason (see \cref{sec:cordic-analysis}).
  \item \textbf{Small normal numbers}, with $|x| \le \pi/4$, already lie in
    the reduced interval, so $k = 0$ and no quadrant correction is
    needed.
  \item \textbf{All other normal numbers} go through range reduction (\cref{sec:range-reduction})
    and the approximation method.
  \item \textbf{Infinity} (exponent 255, fraction 0) produces NaN, since
    $\sin(\pm\infty)$ is an invalid operation~\cite{ieee754}.
  \item \textbf{NaN} inputs (exponent 255, fraction $\neq 0$) produce NaN.
\end{enumerate}

The inputs are therefore divided into two paths: a \emph{bypass path}, for
which the result is known immediately (cases 1--3, 6 and 7), and a
\emph{compute path}, which uses range reduction and the approximation method
(cases 4 and 5).

\section{Range reduction for binary32 inputs}
\label{sec:range-reduction}

Since sine is periodic, every input $x$ is first decomposed into an integer
$k$ and a small remainder $r$~\cite{ly2026payne}:
\begin{equation}
  x = k\,\frac{\pi}{2} + r, \qquad |r| \le \frac{\pi}{4}, \quad k \in \mathbb{Z},
  \label{eq:sin-reduction}
\end{equation}
where:
\begin{itemize}
  \item $k$ is the number of quarter periods contained in $x$; the value
    $k \bmod 4$ determines the quadrant, and therefore which function and
    sign give the final result (\cref{tab:quadrant});
  \item $r$ is the \emph{reduced argument}, which is passed to the
    approximation method (CORDIC or the lookup table).
\end{itemize}

\begin{table}[H]
  \centering
  \caption{Reconstruction of $\sin x$ from the reduced argument $r$}
  \label{tab:quadrant}
  \begin{tabular}{cc}
    \toprule
    $k \bmod 4$ & $\sin x$ \\
    \midrule
    0 & $\phantom{-}\sin r$ \\
    1 & $\phantom{-}\cos r$ \\
    2 & $-\sin r$ \\
    3 & $-\cos r$ \\
    \bottomrule
  \end{tabular}
\end{table}

\paragraph{Use of symmetry.}
Since sine is an odd function, $\sin(-x) = -\sin x$, the reduction is
performed on $|x|$, and the sign bit of the input is applied to the result at
the very end.

\paragraph{Multiplication by $2/\pi$.}
A normalized binary32 input can be written as $x = M \cdot 2^{E}$, where $M$
is the 24-bit integer formed by the implicit leading bit and the fraction
field, and $E = e - 150$ is derived from the stored exponent $e$, where
$150 = 127 + 23$ accounts for the bias and for the 23 fraction bits that are
shifted into the integer $M$.
Computing $x \cdot \frac{2}{\pi}$ therefore reduces to multiplying the integer $M$ by
the bits of $2/\pi$ and shifting the product by $E$ positions. Only a window
of the bits of $2/\pi$ influences the result~\cite{ly2026payne}: the bits
that end up far above the binary point only add multiples of 4 to $k$, and
since only $k \bmod 4$ is needed, they can be ignored.
The bits that end up far below the binary point are smaller than the required precision, so they
can be ignored as well. The constant $2/\pi$ is therefore stored once, with
enough bits for the largest possible exponent, and the exponent $E$ selects
the window of bits that is used for each input.

\paragraph{Cancellation near multiples of $\pi/2$.}
When $x$ is very close to a multiple of $\pi/2$, the reduced argument $r$
becomes very small. A good example is $x = 355$: since
$355/113 \approx \pi$, the input is almost exactly $226 \cdot \frac{\pi}{2}$,
which gives $k = 226$ and $r \approx 0.0000192$ quarter-turns. In binary,
this value of $r$ starts with about 15 leading zeros, which, as in the case
of small CORDIC results, carry no information. The window of $2/\pi$ must
therefore be wide enough that $r$ still has 24 correct significant bits after
these zeros, and this must hold for the worst-case binary32 input, the one
closest to a multiple of $\pi/2$.

\paragraph{Output of range reduction.}
For each input $x$, this stage produces three values: the sign of the input,
which is applied to the final result; the two bits of $k \bmod 4$, which
select between $\pm\sin r$ and $\pm\cos r$ (\cref{tab:quadrant}); and the
reduced argument $r$, expressed in quarter-turns with $|r| \le 0.5$, which is
passed to the approximation method.

\section{CORDIC rotation for sine evaluation}
\label{sec:cordic-analysis}

\paragraph{Mathematical intuition.}
To obtain the sine of an angle $\theta$, we can imagine rotating the unit vector
$(1, 0)$ by $\theta$: the $y$-coordinate of the rotated vector is then equal to
$\sin\theta$, and its $x$-coordinate to $\cos\theta$
(\cref{fig:cordic-rotation}).

\begin{figure}[H]
  \centering
  \includegraphics[scale=0.67]{cordic-rotation}
  \caption{Rotating the unit vector $(1, 0)$ by an angle $\theta$ gives the point
  $(\cos\theta, \sin\theta)$}
  \label{fig:cordic-rotation}
\end{figure}

More generally, the coordinates of a vector $(x, y)$ rotated by $\theta$ are given by the rotation matrix:
\begin{equation}
  \begin{bmatrix} x' \\ y'
  \end{bmatrix}
  =
  \begin{bmatrix}
    \cos\theta & -\sin\theta \\
    \sin\theta & \cos\theta
  \end{bmatrix}
  \begin{bmatrix} x \\ y
  \end{bmatrix}
  \label{eq:rotation}
\end{equation}
Factoring out $\cos\theta$, we obtain:
\begin{equation}
  \begin{bmatrix} x' \\ y'
  \end{bmatrix}
  = \cos\theta
  \begin{bmatrix}
    1 & -\tan\theta \\
    \tan\theta & 1
  \end{bmatrix}
  \begin{bmatrix} x \\ y
  \end{bmatrix}
  \label{eq:rotation-tan}
\end{equation}

The idea of CORDIC is to replace the rotation by $\theta$ with a sequence of smaller rotations by fixed angles $\theta_{i}$, each performed in a direction
$d_{i} = \pm 1$, such that
\begin{equation}
  \theta \approx \sum_{i=0}^{n-1} d_{i}\,\theta_{i}.
  \label{eq:angle-sum}
\end{equation}
The angles are chosen such that $\tan\theta_{i} = 2^{-i}$, that is $\theta_{i} = \arctan 2^{-i}$, so that multiplying by $\tan\theta_{i}$ reduces to a right shift by $i$ positions. Since the same angles are used for every input, the values $\arctan 2^{-i}$ can be precomputed and stored in a table
(\cref{tab:precompval}).

\begin{table}[H]
  \centering
  \caption{Rotation angles $\theta_{i} = \arctan 2^{-i}$ used by CORDIC}
  \label{tab:precompval}
  \begin{tabular}{rcccc}
    \toprule
    $i$ & $\tan\theta_{i}$ & $\theta_{i}$ (deg) & $\theta_{i}$ (rad)
    & $\theta_{i}$ (quarter-turns) \\
    \midrule
    0  & $1$       & 45.00000000 & 0.7853981634 & 0.5000000000 \\
    1  & $2^{-1}$  & 26.56505118 & 0.4636476090 & 0.2951672353 \\
    2  & $2^{-2}$  & 14.03624347 & 0.2449786631 & 0.1559582608 \\
    3  & $2^{-3}$  & 7.12501635  & 0.1243549945 & 0.0791668483 \\
    $\dots$ & $\dots$ & $\dots$ & $\dots$ & $\dots$ \\
    14 & $2^{-14}$ & 0.00349706  & 0.0000610352 & 0.0000388562 \\
    15 & $2^{-15}$ & 0.00174853  & 0.0000305176 & 0.0000194281 \\
    \bottomrule
  \end{tabular}
\end{table}

Applying \cref{eq:rotation-tan} for a single step $i$, with
$\cos(d_{i}\theta_{i}) = \cos\theta_{i}$, gives:
\begin{equation}
  \begin{bmatrix} x_{i+1} \\ y_{i+1}
  \end{bmatrix}
  = \cos\theta_{i}
  \begin{bmatrix}
    1 & -d_{i}\tan\theta_{i} \\
    d_{i}\tan\theta_{i} & 1
  \end{bmatrix}
  \begin{bmatrix} x_{i} \\ y_{i}
  \end{bmatrix}
  \label{eq:rotation-iter}
\end{equation}
After $n$ iterations, we have:
\begin{equation}
  \begin{bmatrix} x_{n} \\ y_{n}
  \end{bmatrix}
  = \left( \prod_{i=0}^{n-1} \cos\theta_{i} \right)
  \begin{bmatrix} 1 & -d_{n-1}\tan\theta_{n-1} \\ d_{n-1}\tan\theta_{n-1} & 1
  \end{bmatrix}
  \cdots
  \begin{bmatrix} 1 & -d_{0}\tan\theta_{0} \\ d_{0}\tan\theta_{0} & 1
  \end{bmatrix}
  \begin{bmatrix} x_{0} \\ y_{0}
  \end{bmatrix}
  \label{eq:rotation-n}
\end{equation}
Since $\tan\theta_{i} = 2^{-i}$, the scale factor becomes
\begin{equation}
  \prod_{i=0}^{n-1} \cos\theta_{i}
  = \prod_{i=0}^{n-1} \frac{1}{\sqrt{1 + 2^{-2i}}}
  \;\xrightarrow{\,n \to \infty\,}\; \frac{1}{K} \approx 0.6073,
  \label{eq:cordic-gain}
\end{equation}
where $K \approx 1.6468$ is the CORDIC gain~\cite{muller2020}. This factor does not depend on the directions $d_{i}$, so it is the same for every input and can be compensated once, by starting from $x_{0} = 1/K$ instead of $1$.

Dropping the scale factor, and introducing a variable $z$ that keeps track of
the angle still left to rotate, gives the CORDIC iteration presented in
\cref{eq:cordic}~\cite{muller2020}, where the values $\arctan 2^{-i}$ are read
from \cref{tab:precompval}.

To compute $\sin \theta$, we use the \emph{rotation} mode of CORDIC, in which $d_{i} = \operatorname{sign}(z_{i})$, so that $z$ is driven towards $0$. Starting from $x_{0} = 1/K$, $y_{0} = 0$ and $z_{0} = \theta$, after $n$ iterations we obtain $y_{n} \approx \sin\theta$ and $x_{n} \approx \cos\theta$.

\paragraph{Convergence range.}
CORDIC reaches the desired angle by adding or subtracting fixed angles:
\ang{45}, then \ang{26.6}, then \ang{14.0}, and so on. Even if all rotations
are performed in the same direction, the largest angle that can be reached is
\begin{equation}
  \sum_{i=0}^{\infty} \arctan 2^{-i} \approx \SI{1.7433}{\radian}
  \approx \ang{99.9},
  \label{eq:cordic-convergence}
\end{equation}
so the iteration converges only for $|z_{0}| \le 1.7433$~rad~\cite{muller2020}.
After range reduction, the reduced argument satisfies $|r| \le \pi/4 \approx
0.785$~rad (\cref{eq:sin-reduction}), which lies well inside this interval.
Therefore, CORDIC converges for every input of the design, and no additional
pre-rotation step is needed.

\paragraph{Angle representation.}
Range reduction (\cref{eq:sin-reduction}) multiplies the input by $2/\pi$, so
its result is naturally expressed in \emph{quarter-turns}, where one unit
corresponds to $\ang{90} = \pi/2$~rad. The integer nearest to this product gives $k$, which selects the quadrant,
while the difference between the product and $k$ gives the reduced argument
$r$, which represents the position within the quadrant and satisfies
$|r| \le 0.5$. Converting $r$ back to radians would require an additional
multiplication by $\pi/2$. Instead, $r$ is kept in quarter-turns, and the
CORDIC rotation angles are stored in the same unit, as
$\theta_{i} / (\pi/2)$ (\cref{tab:precompval}). The iteration in
\cref{eq:cordic} remains unchanged; only the stored constants differ.

\paragraph{Sources of error.}
The accuracy of the fixed-point CORDIC iteration is limited by three sources
of error:
\begin{enumerate}
  \item \textbf{The residual angle.} After $n$ iterations, the angle that
    remains to be rotated, $z_{n}$, is bounded by the last rotation angle,
    $\theta_{n-1} \approx 2^{-(n-1)}$, so each iteration adds roughly one
    correct bit to the result~\cite{muller2020}.
  \item \textbf{Truncation in the shifts.} Each right shift by $i$ positions
    discards the bits that fall off the end of the word. Over many stages
    these losses accumulate, so the datapath needs extra guard bits beyond
    the required output precision.
  \item \textbf{Rounding of the table constants.} Each value
    $\arctan 2^{-i}$ is stored with finite precision, which adds a small
    error at every stage.
\end{enumerate}

\paragraph{Precision of small results.}
\label{pa:precision}
Inside the CORDIC iteration, values are represented in fixed point, so the
position of the binary point never changes. A small result, such as
$\sin r \approx 0.0002$, is therefore represented as
\begin{equation*}
  0.0002 = (0.\underbrace{000000000000}_{12\ \text{zeros}}
  110100011011011100010111\ldots)_{2}.
\end{equation*}
The leading zeros carry no information other than the magnitude of the
number. When the result is converted to the binary32 format, normalization
shifts the significand to the left until the first non-zero bit reaches the
integer position, and the shift is recorded in the exponent:
\begin{equation*}
  0.0002 \approx (1.10100011011011100010111)_{2} \times 2^{-13}.
\end{equation*}
The 12 leading zeros are discarded, and the 24 bits that follow them form the
significand of the result. Since all of these bits must be correct, the
fixed-point computation must be accurate down to bit position
$12 + 24 = 36$, and not only to position 24. The same effect appears in
decimal notation: obtaining three significant digits of $0.000123$ requires
computing six decimal places, since the first three are zeros.

\paragraph{Required number of iterations and word length.}
\label{pa:small-inputs}
For a small reduced argument, $|r| < 2^{-12}$, the approximation
$\sin r \approx r$ is used directly, since its relative error, approximately
$r^{2}/6$, is below the binary32 precision. This covers both small inputs and
inputs close to a multiple of $\pi/2$ (\cref{sec:range-reduction}). The
smallest argument that reaches the CORDIC iteration is therefore about
$2^{-12}$, which corresponds to the case illustrated above: the result has
about 12 leading zeros, so an absolute accuracy of approximately $2^{-36}$ is
required. Since each iteration adds roughly one correct bit, this requires
about 36 to 38 iterations; with a small safety margin, the design uses about
40 iterations. In addition, the truncation errors of all stages accumulate,
so a few guard bits must be added beyond the 36 bits required by the result,
which leads to an internal word length of about 44 fractional bits.

\section{Lookup table with linear interpolation}
\label{sec:lut}

The second approximation method evaluates $\sin r$ from a precomputed table.
This section determines how the table must be built and used to reach
binary32 accuracy.

\paragraph{Limits of direct tabulation.}
A plain table stores the values of the sine function at $N$ evenly spaced
points over one quadrant, $r \in [0, 1]$ quarter-turns, and every input that
lies between two consecutive points receives the value of the point below it.
The number of correct bits of this approximation grows only logarithmically
with $N$: doubling the size of the table adds just one correct bit, even when
the stored values are chosen optimally (\cref{eq:tab-optimal})~\cite{muller2020}.
A table with 4096 entries reaches only about 11 correct bits, and a 24-bit
result would require around 30 million entries.

\paragraph{Linear interpolation.}
Instead of using only the lower neighbour, both neighbours of the input can be
used, and the sine function can be approximated by the straight line between
them. If the input lies between the table points $i$ and $i+1$, at a fraction
$d$ of the distance between them ($0 \le d < 1$), then
\begin{equation}
  \sin r \approx T[i] + d \cdot \big( T[i+1] - T[i] \big),
  \label{eq:lerp}
\end{equation}
where $T[i]$ is the stored value at point $i$. Each evaluation therefore
requires two table reads, one multiplication and two additions.

\paragraph{Required table size.}
The error of linear interpolation is no longer proportional to the distance
$h$ between table points, but to $h^{2}$, since the error of a first-order
approximation grows with the square of the distance to the reference
point~\cite{muller2020}. Each doubling of the table now
adds two correct bits instead of one. \Cref{tab:lut-size} compares the two
approaches, as obtained by evaluating both methods numerically over the whole
quadrant. With interpolation, 4096 entries are sufficient for about 25.7
correct bits, which exceeds the binary32 precision, while a plain table would
need around 30 million entries. For the plain table, the absolute error is
given, since its relative error near zero is unbounded. For linear interpolation, the errors in
\cref{tab:lut-size} are relative, and they remain small even close to zero: since the table starts
exactly at $\sin 0 = 0$, the line and the curve shrink together near the
origin. Arguments smaller than $2^{-12}$ are handled directly, as described in
\cref{sec:cordic-analysis}.

\begin{table}[H]
  \centering
  \caption{Correct bits of the table-based approximation, with and without
  linear interpolation}
  \label{tab:lut-size}
  \begin{tabular}{rcc}
    \toprule
    Table entries $N$ & Plain table & Linear interpolation \\
    \midrule
    256  & $\approx 7$  & $\approx 17$   \\
    1024 & $\approx 9$  & $\approx 21$   \\
    2048 & $\approx 10$ & $\approx 23$   \\
    4096 & $\approx 11$ & $\approx 25.7$ \\
    \bottomrule
  \end{tabular}
\end{table}

\paragraph{One table for both sine and cosine.}
As shown in \cref{tab:quadrant}, range reduction requires either $\sin r$ or
$\cos r$, depending on the quadrant. The CORDIC iteration provides both values
at the same time, whereas a table-based method would in principle need two
separate tables. However, with angles expressed in quarter-turns,
\begin{equation}
  \cos\!\left(r \cdot \frac{\pi}{2}\right)
  = \sin\!\left((1 - r) \cdot \frac{\pi}{2}\right),
  \label{eq:sincosqt}
\end{equation}
so the cosine is simply the sine read from the other end of the quadrant. A
single table of the sine function over one quadrant therefore covers both
cases: it is read at $|r|$ when the sine is needed, and at $1 - |r|$ when the
cosine is needed. Since the argument can reach the end of the quadrant, the
table stores $N + 1 = 4097$ values, including the endpoint
$\sin(\pi/2) = 1$, which is also required for interpolating the last interval.

\paragraph{Table addressing.}
With $4096 = 2^{12}$ entries, the argument of the table, expressed in
quarter-turns and lying in the interval $[0, 1]$, is divided directly by its
bits:
\begin{equation*}
  (0.\underbrace{101100111010}_{\text{index } i}\,
  \underbrace{01101001\ldots}_{\text{fraction } d})_{2}.
\end{equation*}
The 12 most significant fractional bits select the table entry $i$, while
the remaining bits form the fraction $d$, which gives the position of the
argument between the entries $i$ and $i+1$ (\cref{eq:lerp}). Addressing the
table therefore requires no arithmetic, only the selection of two bit fields.
This is possible because the argument is expressed in quarter-turns: one
quadrant corresponds exactly to the interval $[0, 1]$, so its binary
fraction can be used as a table address directly.

\section{Normalization and rounding to binary32}
\label{sec:normalization}

The outputs of both CORDIC and the table-based method are fixed-point numbers
between 0 and 1, with about 44 fractional bits. Range reduction also provides
the sign of the input and $k \bmod 4$.
This stage will turn this number into a binary32 format in four steps.

\paragraph{Sign of the result.}
Two values determine the sign of the result: the sign of the input and the
quadrant given by $k \bmod 4$. For $k \bmod 4 = 2$ or $3$, i.e.\ in the third
and fourth quadrants, the result is negative (\cref{tab:quadrant}); this
corresponds to the most significant bit of $k \bmod 4$ being 1. The sign of
the result is therefore obtained as the exclusive OR of the input sign and
this bit, since two negative factors cancel each other.
Since the approximation method operates on $|r|$, the sign of $r$ is also
carried as side information; it changes the sign of the result only when the
sine is selected ($k \bmod 4 = 0$ or $2$), since $\sin(-r) = -\sin r$ while
$\cos(-r) = \cos r$.

\paragraph{Leading-zero count.}
The fixed-point result produced by the approximation method lies in the
interval $[0, 1]$, and small results are preceded by a number of leading
zeros. The next step is to count these zeros, which determines the position
of the first non-zero bit and therefore the magnitude of the result.
Let $z$ denote the number of leading zeros; the first non-zero bit is then
at position $z + 1$.

\paragraph{Normalization and exponent.}
The fixed-point value is then shifted to the left by $z$ positions, so that
the first non-zero bit reaches the position in front of the binary point.
This bit becomes the implicit leading bit of the binary32 format, and only
the bits that follow it are kept as the fraction. Since the value has the
form $1.f \times 2^{-(z+1)}$, the stored exponent, including the bias of 127,
is
\begin{equation}
  e = 127 - (z + 1) = 126 - z.
  \label{eq:norm-exponent}
\end{equation}

\paragraph{Rounding to 24 bits.}
The significand of a binary32 number has only 24 bits: the implicit leading
bit and the 23 fraction bits. The bits of the fixed-point result that follow
these 24 bits must therefore be discarded, and the result is rounded according
to the default rounding mode of the standard, \emph{round to nearest, ties to
even}~\cite{ieee754}:
\begin{itemize}
  \item if the first discarded bit is 0, the result is rounded down, i.e.\ the
    discarded bits are simply removed;
  \item if the first discarded bit is 1 and at least one of the following bits
    is 1, the result is rounded up, by adding one unit to the last kept
    bit;
  \item if the first discarded bit is 1 and all following bits are 0, the
    value lies exactly halfway between two candidates, and the one whose
    last bit is 0 is chosen, which avoids a systematic rounding bias.
\end{itemize}
Rounding up can cause an overflow of the significand: a value whose 24 bits
are all 1, $(1.11\ldots1)_{2}$, becomes $(10.00\ldots0)_{2}$ after one unit
is added to its last bit. In this case, the significand is shifted right by
one position and the exponent $e$ (\cref{eq:norm-exponent}) is increased by
one. This occurs, for instance, for results very close to 1.

It is worth mentioning that the value being rounded is already an approximation,
so the result is not guaranteed to be correctly rounded. It can rarely differ
by one unit in the last place, when the true value lies very close to a rounding boundary.

Bypassed inputs ($x$ itself, $\pm 0$, NaN) skip this stage,
except the tiny-$r$ case, which goes through normalization.

\section{Pipelined execution}
\label{sec:pipeline}

Both approximation methods are organized as a pipeline, which accepts a new
input at every clock cycle. CORDIC is particularly well suited to this
organization: since its iterations all have the same structure, each
iteration can be assigned to its own pipeline stage.

\paragraph{Iterative vs pipelined execution.}
An iterative algorithm with $n$ iterations can be executed in two ways:
\emph{iteratively}, by a single unit that gets reused for each iteration,
or in a \emph{pipeline}, by $n$ (mostly) identical units,
each performing one iteration and separated by pipeline registers.
\Cref{tab:iter-vs-pipe} compares the two variants.

\begin{table}[H]
  \centering
  \caption{Iterative and pipelined execution of an algorithm with $n$
  iterations}
  \label{tab:iter-vs-pipe}
  \begin{tabular}{lcc}
    \toprule
    & Iterative & Pipelined \\
    \midrule
    Iteration units            & 1                & $n$             \\
    Latency                    & $n$ cycles       & $n$ cycles      \\
    New input accepted         & every $n$ cycles & every cycle     \\
    Inputs processed at once   & 1                & up to $n$       \\
    \bottomrule
  \end{tabular}
\end{table}

The latency is the same in both cases, since every input still needs $n$ iterations,
but the throughput of the pipeline is $n$ times higher, and a new result is produced
at each clock cycle. However, for the pipeline the price grows proportionally to the
number of pipeline stages.

\paragraph{Structure of one pipeline stage.}
In a pipelined form, stage $i$ always does iteration $i$ of \cref{eq:cordic}, so
the shift amount is fixed (stage 5 always shifts by 5)
and the angle constant is fixed (stage 5 always uses $\arctan 2^{-5}$). In the iterative form,
by contrast, the single unit is reused for every iteration, so it must shift by a different
amount and use a different constant in each cycle, which requires a variable shifter and a
table of constants.
Each stage consists only of three additions or subtractions, plus a sign check of $z_{i}$.
The fixed shifts require no computation, only a selection of the bits.

\paragraph{Side information.}
\label{pa:side-info}
Each input carries extra information through the pipeline: the input sign, $k \bmod 4$,
and the type of input (\cref{sec:input-classification}). This information advances
one stage per clock cycle
together with the computed data so they arrive at the output at the same time.
Bypassed inputs also must wait for the same number of cycles, otherwise the results
would come out of order.

\paragraph{The slowest stage sets the pace.}
All stages share the same clock, so the clock can only be as fast
as the slowest stage. The largest single operation of the design is the multiplication of the 24-bit significand
with the selected window of $2/\pi$ in range reduction (\cref{sec:range-reduction}).
If this operation cannot be completed within one clock cycle, it is divided
over several pipeline stages, at the cost of a slightly higher latency.

\paragraph{Analysing latency.}
The CORDIC method has around 40 approximation stages (one per iteration), with 3 additions per stage and no
multiplication, meaning that the total latency is relatively high --- every input reaches the output after
around 50 clock cycles, which include the range reduction, input classification and normalization after leaving
the approximation stages. The cost of the approximation part is $40 \times 3$ adders.

The table-based method requires only a few approximation stages, which follow \cref{eq:lerp}: reading the
entries $T[i]$ and $T[i+1]$, computing their difference, multiplying it by $d$ and adding the result to
$T[i]$. Including the same preparation and normalization stages, its latency
is around 10--12 clock cycles, and its approximation part requires one table
of 4097 values and one multiplier.

The two chosen methods represent opposite trade-offs: CORDIC avoids
multiplications and large memories at the cost of a longer pipeline, while
the table-based method has a shorter pipeline, but requires a memory and a
multiplier. Measuring the costs on the same device is one of the main objectives of this project
(\cref{sec:objectives}).

\section{Flow control with backpressure}
\label{sec:backpressure}

The pipeline produces one result per clock cycle, but the receiver may
refuse a result at any time through the valid/ready handshake
(\cref{rule:valid-hold}). Since up to one input per stage is in flight, the
pipeline must react without losing data. Each stage carries a valid bit that
marks whether it holds a real input. Two strategies are possible: stalling
all stages while the receiver is not ready, which needs no extra storage but
requires a control signal that reaches every stage in the same cycle; or
letting the pipeline run and collecting results in an output buffer, whose
size must match the pipeline depth.